\documentclass{article}
\newcommand{\pair}[2]{(#1, \, #2)}
\newcommand{\ineq}{%
  \mathrel{\ooalign{%
    $\in$\cr
    \hidewidth\raisebox{-0.6ex}{\rule{0.7em}{0.6pt}}\hidewidth\cr
  }}}

\begin{document}
  Antes de iniciar, recalcalcar que la definición de natural se toma bajo el supuesto:

  \begin{itemize}
    \item $x$ es transitivo.
    \item $\pair{x}{\in}$ es un orden estricto y total  en $x$.
    \item Todo subconjunto no vacío de $x$ tiene máximo y mínimo.
  \end{itemize}

  También debemos considerar el axioma del ínfinito para asegurar que existe al menos 
  un conjunto transitivo.
  \newline\newline


  \textit{Teorema}: Si $x$ es natural y $A$ es un conjunto inductivo, $x \in A$.
  \newline\newline

  Por ser $A$ inductivo sabemos que $\emptyset \in A$. Supongamos que existe 
  un natural $x$ tal que $x \not\in A$ y tomemos $x\setminus A$. Si este conjunto es 
  vacío, es decir, $x\setminus A = \emptyset$, se puede decir que $A \subseteq x$
  y como $\emptyset \in A$, $A$ es un subconjunto no vacío de $x$, entonces 
  $A$ debería tener un máximo $m \in A$, pero como $A$ es inductivo, $m^+ \in A$, 
  contradiciendo así que $m$ sea máximo.
  \newline\newline


  Con esto último, se concluye que $x \setminus A \neq \emptyset$. Este es un subconjunto de $x$ no vacío,
  entonces tiene mínimo, denotemoslo $y$. Es decir, si $z \in y$, entonces $z \in A$.
  De esto, $y \neq \emptyset$. Llamamos $d$ el máximo elemento menor que $y$ (¿Por qué se puede esto?)
  y sabemos inmediatamente que $d \in A$.
  \newline\newline


  Como $y$ es natural, entonces $d \subseteq y$ y $\{d\} \subseteq y$, entonces $d^+ \subseteq y$.
  Si $z \in y$, entonces $z \ineq d$, luego $z \in d^+$, concluyendo así que 
  $d^+ = y$, pero como $d \in A$, entonces $y = d^+ \in A$, lo cual es una clara 
  contradicción. El único supuesto que hicimos es que $x \not\in A$, luego concluimos
  que $x \in A$.
\end{document}
